%------------------------------------------------------------------------------%

\usepackage{integracao_e_series}

\title {Integrais Indefinidas}

%------------------------------------------------------------------------------%
\begin{document}

\addtitle

%------------------------------------------------------------------------------%
\section{Integrais Indefinidas}

%------------------------------------------------------------------------------%
\begin{frame}{Integrais Indefinidas}
  Sinônimos
  \begin{itemize}
    \item Integral indefinida
    \item Primitiva
    \item Antiderivada
  \end{itemize}
\end{frame}

%------------------------------------------------------------------------------%
\begin{frame}{Definição}
  Uma função $F$ é a \emph{Primitiva}
  da função $f$ em um intervalo $I$ se
  \[
    \frac{dF}{dx}(x) = f(x)
  \]
  para todo $x\in I$
\end{frame}

%------------------------------------------------------------------------------%
\begin{frame}{Primitiva Mais Geral}
  Se $G$ é uma primitiva de $f$ no intervalo $I$, então a
  \emph{primitiva mais geral} de $f$ em $I$ é
  \[
    F(x) = G(x) + C
  \]
  onde $C$ é uma constante arbitrária
\end{frame}

%------------------------------------------------------------------------------%
\begin{frame}{Notação}
  Notação da Primitiva, Antiderivada e Integral Indefinida de uma função $f$
  \[
    \int f(x) \, dx
  \]
\end{frame}

%------------------------------------------------------------------------------%
\section{Cálculos e Propriedades}

%------------------------------------------------------------------------------%
\begin{frame}{Calculando Raízes Quadradas}
\(\sqrt{x}\) é a operação inversa de \(x^2\) para \(x\geq0\)

\pause
\vspace{\baselineskip}
\begin{columns}
\column{0.5\linewidth}
  Sabemos calcular o quadrado
  \begin{align*}
    \vphantom{\sqrt{9}} 1^2 & = 1 \\
    \vphantom{\sqrt{9}} 2^2 & = 4 \\
    \vphantom{\sqrt{9}} 3^2 & = 9 \\
    \vphantom{\sqrt{9}} 4^2 & = 16
  \end{align*}
\column{0.5\linewidth}
  \pause
  Valores conhecidos para raízes
  \begin{align*}
    \sqrt{1}  & = 1 \\
    \sqrt{4}  & = 2 \\
    \sqrt{9}  & = 3 \\
    \sqrt{16} & = 4
  \end{align*}
\end{columns}

\pause
\vspace{\baselineskip}
Para os demais valores precisamos de uma nova técnica
\end{frame}

%------------------------------------------------------------------------------%
\begin{frame}{Exemplo}
  \emph{Encontrar a primitiva mais geral de \(f(x)= \cos(x)\)}

  \pause
  \vspace{\baselineskip}
  Sabemos que
  \(
    G(x) = \sin(x)
  \)
  é uma \emph{primitiva} para $f$, pois
  \[
    G'(x)
    \;=\; \frac{d}{dx}\sin(x)
    \;=\; \cos(x) = f(x)
  \]
  \pause
  Portanto, a \emph{primitiva mais geral} é
  \[
    F(x)
    \;=\; \int \cos(x)\,dx
    \;=\; G(x) + C
    \;=\; \sin(x) + C
  \]
\end{frame}

%------------------------------------------------------------------------------%
\begin{frame}{Tabela de Primitivas}
\renewcommand*{\arraystretch}{1.8}
\begin{columns}
\column{0.5\linewidth}
\centering
  \begin{tabular}{c<{\;}>{\;}c}
    \hline
    Função                      & Primitiva               \\ \hline
    \(x^n\), \(n \neq -1\)      & \(\frac{x^{n+1}}{n+1}\) \\[1mm]
    \(\frac{1}{x}\)             & \(\ln\abs{x}\)          \\
    \(e^x\)                     & \(e^x\)                 \\
    \(\sin(x) \)                & \(-\cos(x)\)            \\
    \( \cos(x)\)                & \(\sin(x)\)             \\\hline
  \end{tabular}
\column{0.5\linewidth}
\centering
  \begin{tabular}{c<{\;}>{\;}c}
    \hline
    Função                      & Primitiva               \\ \hline
    \(\sec^2(x)\)               & \(\tan(x)\)             \\
    \(\sec(x) \tan(x)\)         & \(\sec(x)\)             \\[1mm]
    \(\frac{1}{\sqrt{1-x^2}}\)  & \(\arcsin(x)\)          \\[2mm]
    \(\frac{-1}{\sqrt{1-x^2}}\) & \(\arccos(x)\)          \\[2mm]
    \(\frac{1}{1+x^2}\)         & \(\arctan(x)\)          \\[3mm] \hline
  \end{tabular}
\end{columns}
\end{frame}

%------------------------------------------------------------------------------%
\begin{frame}{Linearidade da Integral Indefinida}
  Sejam $f$ e $g$ integráveis e $k$ uma constante temos \\[7mm]
  \pause
  \begin{enumerate}[<+->]
    \item \(\int f(x) \,{\color{blue}+}\, g(x) \,dx
          = \int f(x) \,dx \,{\color{blue}+} \int g(x) \,dx\) \\[7mm]
    \item \(\int {\color{blue}k} f(x) \,dx = {\color{blue}k} \int f(x) \,dx\)
  \end{enumerate}
\end{frame}

%------------------------------------------------------------------------------%
\begin{frame}
  \tableofcontents
\end{frame}

%------------------------------------------------------------------------------%
\end{document}
%------------------------------------------------------------------------------%
